\documentclass[12pt,a4paper]{article}
\usepackage[utf8]{inputenc}
\usepackage[T1]{fontenc}
\usepackage[brazil]{babel}
\usepackage{geometry}
\usepackage{booktabs}
\usepackage{xcolor}
\usepackage{listings}
\usepackage{graphicx}
\usepackage{enumitem}
\usepackage{float}
\usepackage{array}
\usepackage{tabularx}

\geometry{landscape, margin=2.2cm}
\setlength{\parindent}{0pt}
\setlength{\parskip}{0.45em}
\renewcommand{\arraystretch}{1.15}

\title{Requisitos de Verificação -- Bloco ReLU}
\author{Matheus Grossi}
\date{2026}

% Inclui a imagem quando ela estiver disponível no mesmo diretório do .tex.
% Enquanto cap1/cap2 não estiverem presentes, um quadro informativo é exibido.
\newcommand{\coverageimage}[2]{%
    \IfFileExists{#1}{%
        \fbox{\includegraphics[width=0.94\textwidth,height=0.34\textheight,keepaspectratio]{#1}}%
    }{%
        \fbox{\parbox[c][0.24\textheight][c]{0.88\textwidth}{\centering
        \textbf{Imagem de cobertura não disponível no pacote recebido.}\\[0.5em]
        Arquivo esperado: \texttt{#1}}}%
    }%
    \caption{#2}%
}

\begin{document}

\maketitle

\vspace{1.5cm}
\section{Identificação do Bloco}
\begin{itemize}[label=\textbullet]
    \item \textbf{1.1. Nome do bloco:} Unidade de ativação ReLU (Rectified Linear Unit), com largura de 32 bits.
    \item \textbf{1.2. Versão:} 1.0.
    \item \textbf{1.3. Autor do bloco:} Felipe Guedes.
    \item \textbf{1.4. Responsável pela verificação:} Matheus Grossi.
    \item \textbf{1.5. Escopo da verificação:} comportamento funcional da ReLU para entradas positivas, negativas, zero e padrões representativos contendo estados de quatro níveis lógicos.
\end{itemize}

\section{Objetivo da Verificação}
O objetivo da verificação é garantir que o módulo \textbf{ReLU}, configurado para uma largura de 32 bits, apresente o comportamento esperado para valores positivos, negativos e zero.

Como seria inviável testar exaustivamente todas as combinações possíveis de uma entrada de 32 bits, a estratégia de verificação divide o espaço de entrada em rotinas representativas. O comportamento funcional depende principalmente do bit mais significativo (MSB), que indica o sinal do número em representação de complemento de dois.

\begin{itemize}
    \item \texttt{MSB = 0}: a entrada representa um valor não negativo e deve ser propagada integralmente para a saída;
    \item \texttt{MSB = 1}: a entrada representa um valor negativo e a saída deve ser igual a zero.
\end{itemize}

Os 31 bits menos significativos são exercitados com padrões distintos para verificar a propagação dos dados e a decisão funcional associada ao MSB.

\section{Estratégia de Verificação}
Para cada padrão aplicado aos bits \texttt{[30:0]}, os dois estados funcionais do MSB são exercitados. A estratégia é composta pelas seguintes rotinas:

\begin{enumerate}
    \item Bits inferiores preenchidos com valores aleatórios;
    \item Bits inferiores fixados em \texttt{0};
    \item Bits inferiores fixados em \texttt{1};
    \item Bits inferiores alternados entre \texttt{0} e \texttt{1};
    \item Bits inferiores fixados em alta impedância (\texttt{Z});
    \item Bits inferiores fixados em estado indefinido (\texttt{X});
    \item Bits inferiores alternados entre \texttt{X} e \texttt{Z}.
\end{enumerate}

Para os padrões alternados, devem ser exercitadas as duas polaridades da sequência, como \texttt{0101...} e \texttt{1010...}.

\section{Casos de Teste Executados / Planejados}
\begin{table}[H]
\centering
\scriptsize
\begin{tabularx}{\textwidth}{|c|c|p{4.0cm}|X|}
\hline
\textbf{Caso} & \textbf{MSB} & \textbf{Bits [30:0]} & \textbf{Resultado esperado} \\ \hline
1  & 0 & Aleatórios                  & \texttt{out = in} \\ \hline
2  & 1 & Aleatórios                  & \texttt{out = 0} \\ \hline
3  & 0 & Todos \texttt{0}            & \texttt{out = 0} \\ \hline
4  & 1 & Todos \texttt{0}            & \texttt{out = 0} \\ \hline
5  & 0 & Todos \texttt{1}            & \texttt{out = in} \\ \hline
6  & 1 & Todos \texttt{1}            & \texttt{out = 0} \\ \hline
7  & 0 & Alternados \texttt{0/1}     & \texttt{out = in} \\ \hline
8  & 1 & Alternados \texttt{0/1}     & \texttt{out = 0} \\ \hline
9  & 0 & Todos \texttt{Z}            & \texttt{out = in}, preservando os estados \texttt{Z} dos bits inferiores \\ \hline
10 & 1 & Todos \texttt{Z}            & \texttt{out = 0} \\ \hline
11 & 0 & Todos \texttt{X}            & \texttt{out = in}, preservando os estados indefinidos dos bits inferiores \\ \hline
12 & 1 & Todos \texttt{X}            & \texttt{out = 0} \\ \hline
13 & 0 & Alternados entre \texttt{X/Z} & \texttt{out = in}, preservando os estados de quatro níveis aplicados à entrada \\ \hline
14 & 1 & Alternados entre \texttt{X/Z} & \texttt{out = 0} \\ \hline
\end{tabularx}
\caption{Matriz principal de verificação da ReLU.}
\end{table}

\section{Casos de Fronteira}
Os seguintes valores possuem importância especial e devem ser garantidos explicitamente durante a verificação:

\begin{table}[H]
\centering
\small
\begin{tabular}{|p{7cm}|p{5cm}|p{6cm}|}
\hline
\textbf{Condição} & \textbf{Valor representado} & \textbf{Resultado esperado} \\ \hline
Todos os bits em \texttt{0} & 0 & 0 \\ \hline
MSB \texttt{0} e demais bits \texttt{1} & Maior positivo representável & Própria entrada \\ \hline
MSB \texttt{1} e demais bits \texttt{0} & Menor negativo representável & 0 \\ \hline
Todos os bits em \texttt{1} & -1 & 0 \\ \hline
\end{tabular}
\caption{Casos de fronteira da função ReLU para 32 bits.}
\end{table}

Esses quatro casos verificam os limites das duas regiões de funcionamento da ReLU: propagação para números não negativos e saturação em zero para números negativos.

\section{Resultados}
\subsection{Comportamentos esperados}
A verificação deve confirmar que valores com MSB igual a \texttt{0} são propagados para a saída e que valores com MSB igual a \texttt{1} produzem saída nula. O valor zero também deve permanecer em zero.

Para os padrões contendo \texttt{X} ou \texttt{Z} nos 31 bits inferiores, o objetivo é verificar que esses estados não alteram a decisão determinada pelo MSB conhecido. Quando o MSB é \texttt{0}, os bits inferiores são propagados; quando o MSB é \texttt{1}, a saída é forçada a zero.

\subsection{Critério de aprovação}
O bloco será considerado aprovado quando:
\begin{itemize}
    \item todos os casos determinísticos apresentarem o resultado esperado;
    \item todos os estímulos aleatórios apresentarem o resultado esperado;
    \item valores positivos forem preservados integralmente;
    \item valores negativos produzirem saída igual a zero;
    \item o valor zero produzir saída igual a zero;
    \item estados \texttt{X} e \texttt{Z} nos bits inferiores não alterarem a decisão determinada por um MSB conhecido;
    \item todas as classes definidas na cobertura funcional forem exercitadas;
    \item nenhum erro de comparação for identificado durante a simulação.
\end{itemize}

\section{Análise dos Resultados}
\subsection{Interpretação dos casos de sucesso/falha}
\textbf{Sucesso:} um caso é considerado correto quando a saída observada corresponde à função ReLU definida para o valor do MSB e preserva, quando aplicável, o conteúdo dos bits inferiores.\par

\textbf{Falha:} ocorre quando um valor não negativo não é propagado corretamente, quando uma entrada negativa não resulta em zero ou quando os padrões especiais aplicados aos bits inferiores produzem comportamento incompatível com a decisão estabelecida pelo MSB.

\subsection{Sumário da cobertura de verificação}
A cobertura deve garantir o exercício das combinações entre os dois estados funcionais do MSB e os padrões representativos definidos para os demais bits. O plano original estabelece como cobertura funcional mínima o cruzamento do MSB com cinco classes principais de padrão: aleatório, todos em zero, todos em um, alternados e todos em \texttt{Z}. Além disso, a matriz de testes também inclui estímulos com \texttt{X} e alternância \texttt{X/Z}.

\begin{figure}[H]
    \centering
    \coverageimage{cap1.png}{Cobertura da verificação da ReLU -- captura 1.}
\end{figure}

\begin{figure}[H]
    \centering
    \coverageimage{cap2.png}{Cobertura da verificação da ReLU -- captura 2.}
\end{figure}

As duas capturas acima devem ser usadas como evidência da taxa de cobertura obtida na execução da bancada. A interpretação quantitativa dos percentuais deve ser feita diretamente a partir dos valores apresentados nas imagens, evitando extrapolar métricas que não estejam explicitamente registradas no relatório de cobertura.

\subsection{Comentários sobre limitações dos testes}
A verificação não busca percorrer todas as $2^{32}$ combinações binárias possíveis. Em vez disso, o espaço de entrada é particionado segundo o comportamento funcional da ReLU e exercitado com padrões representativos. Essa abordagem reduz o volume de simulação sem deixar de testar diretamente a decisão de sinal, os principais casos de fronteira e padrões especiais de quatro níveis lógicos.

\section{Conclusão}
O plano de verificação da ReLU foi estruturado para validar de forma direcionada a regra funcional do bloco: preservar entradas não negativas e produzir zero para entradas negativas. A estratégia utiliza o MSB como principal eixo de decisão e combina seus dois estados com diferentes padrões nos 31 bits inferiores, incluindo valores aleatórios, constantes, alternados, \texttt{X} e \texttt{Z}.

O fechamento da verificação deve considerar simultaneamente os resultados das comparações funcionais e as métricas apresentadas nas capturas de cobertura \texttt{cap1} e \texttt{cap2}.

\vfill
\begin{center}
    \textbf{Bloco verificado:} ReLU -- 32 bits
\end{center}

\end{document}
